% GENERATED FILE. Edit canonical lessons, then run npm run generate:books.
% canonical-source-hash: fnv1a64:2ed942fb9f070bf9
% canonical-lessons: SA-C169-parinamah, SA-C169-ankah, SA-C169-sreni, SA-C169-sabdakosah, SA-C169-lekhanam

\chapter{Result, Mark, Row, Dictionary, Writing}
\label{ch:sa-result-mark-row-dictionary-writing}
\hlchaptermodality{169}

\begin{chapteropening}
\textbf{By the end of this chapter:} \emph{I can say a result, a mark, a row, a dictionary and writing.}

Complete the last lesson of chapter 169: I can say a result, a mark, a row, a dictionary and writing.
\end{chapteropening}

\section[\texorpdfstring{\sk{परिणामः} (pariṇāmaḥ)}{परिणामः (pariṇāmaḥ)}]{\sk{परिणामः} (pariṇāmaḥ) — a result}

\label{lesson:SA-C169-parinamah}



\begin{quote}
Before the new one: say the Sanskrit for a colour, then the Sanskrit for a word.
\end{quote}

\subsection*{You'll want to know: \sk{परिणामः}}
\textbf{\sk{परिणामः}} — \emph{pariṇāmaḥ} — \textquotedblleft{}a result\textquotedblright{}.

\begin{grammarlens}[title={What you've built}]
One more word for this chapter.
\end{grammarlens}

\subsection*{Your turn}
\begin{itemize}
\raggedright
  \item \emph{Say it:} \emph{pariṇāmaḥ}
  \item \emph{Say it:} \emph{pariṇāmaḥ}, once more
  \item \emph{Say it:} say \emph{pariṇāmaḥ}
  \item \emph{Recall:} say the Sanskrit for a colour, then the Sanskrit for a word, then say \emph{pariṇāmaḥ} again
\end{itemize}

\begin{tcolorbox}[breakable,colback=teal!4,colframe=teal!35!black,title={Before you move on}]
What does \sk{परिणामः} mean? (\textquotedblleft{}A result\textquotedblright{}.) Say it once more.
\end{tcolorbox}
\section[\texorpdfstring{\sk{अंकः} (aṅkaḥ)}{अंकः (aṅkaḥ)}]{\sk{अंकः} (aṅkaḥ) — a mark, a number}

\label{lesson:SA-C169-ankah}



\begin{quote}
Before the new one: say the Sanskrit for a word, then the Sanskrit for a result.
\end{quote}

\subsection*{You'll want to know: \sk{अंकः}}
\textbf{\sk{अंकः}} — \emph{aṅkaḥ} — \textquotedblleft{}a mark, a number\textquotedblright{}.

\begin{grammarlens}[title={What you've built}]
One more word for this chapter.
\end{grammarlens}

\subsection*{Your turn}
\begin{itemize}
\raggedright
  \item \emph{Say it:} \emph{aṅkaḥ}
  \item \emph{Say it:} \emph{aṅkaḥ}, once more
  \item \emph{Say it:} say \emph{aṅkaḥ}
  \item \emph{Recall:} say the Sanskrit for a word, then the Sanskrit for a result, then say \emph{aṅkaḥ} again
\end{itemize}

\begin{tcolorbox}[breakable,colback=teal!4,colframe=teal!35!black,title={Before you move on}]
What does \sk{अंकः} mean? (\textquotedblleft{}A mark, a number\textquotedblright{}.) Say it once more.
\end{tcolorbox}
\section[\texorpdfstring{\sk{श्रेणी} (śreṇī)}{श्रेणी (śreṇī)}]{\sk{श्रेणी} (śreṇī) — a row, a grade}

\label{lesson:SA-C169-sreni}



\begin{quote}
Before the new one: say the Sanskrit for a result, then the Sanskrit for a mark.
\end{quote}

\subsection*{You'll want to know: \sk{श्रेणी}}
\textbf{\sk{श्रेणी}} — \emph{śreṇī} — \textquotedblleft{}a row, a grade\textquotedblright{}.

\begin{grammarlens}[title={What you've built}]
One more word for this chapter.
\end{grammarlens}

\subsection*{Your turn}
\begin{itemize}
\raggedright
  \item \emph{Say it:} \emph{śreṇī}
  \item \emph{Say it:} \emph{śreṇī}, once more
  \item \emph{Say it:} say \emph{śreṇī}
  \item \emph{Recall:} say the Sanskrit for a result, then the Sanskrit for a mark, then say \emph{śreṇī} again
\end{itemize}

\begin{tcolorbox}[breakable,colback=teal!4,colframe=teal!35!black,title={Before you move on}]
What does \sk{श्रेणी} mean? (\textquotedblleft{}A row, a grade\textquotedblright{}.) Say it once more.
\end{tcolorbox}
\section[\texorpdfstring{\sk{शब्दकोशः} (śabdakośaḥ)}{शब्दकोशः (śabdakośaḥ)}]{\sk{शब्दकोशः} (śabdakośaḥ) — a dictionary}

\label{lesson:SA-C169-sabdakosah}



\begin{quote}
Before the new one: say the Sanskrit for a mark, then the Sanskrit for a row.
\end{quote}

\subsection*{You'll want to know: \sk{शब्दकोशः}}
\textbf{\sk{शब्दकोशः}} — \emph{śabdakośaḥ} — \textquotedblleft{}a dictionary\textquotedblright{}.

\begin{grammarlens}[title={What you've built}]
One more word for this chapter.
\end{grammarlens}

\subsection*{Your turn}
\begin{itemize}
\raggedright
  \item \emph{Say it:} \emph{śabdakośaḥ}
  \item \emph{Say it:} \emph{śabdakośaḥ}, once more
  \item \emph{Say it:} say \emph{śabdakośaḥ}
  \item \emph{Recall:} say the Sanskrit for a mark, then the Sanskrit for a row, then say \emph{śabdakośaḥ} again
\end{itemize}

\begin{tcolorbox}[breakable,colback=teal!4,colframe=teal!35!black,title={Before you move on}]
What does \sk{शब्दकोशः} mean? (\textquotedblleft{}A dictionary\textquotedblright{}.) Say it once more.
\end{tcolorbox}
\section[\texorpdfstring{\sk{लेखनम्} (lekhanam)}{लेखनम् (lekhanam)}]{\sk{लेखनम्} (lekhanam) — writing}

\label{lesson:SA-C169-lekhanam}



\begin{quote}
Before the new one: say the Sanskrit for a row, then the Sanskrit for a dictionary.
\end{quote}

\subsection*{You'll want to know: \sk{लेखनम्}}
\textbf{\sk{लेखनम्}} — \emph{lekhanam} — \textquotedblleft{}writing\textquotedblright{}.

\begin{grammarlens}[title={What you've built}]
One more word for this chapter.
\end{grammarlens}

\subsection*{Your turn}
\begin{itemize}
\raggedright
  \item \emph{Say it:} \emph{lekhanam}
  \item \emph{Say it:} \emph{lekhanam}, once more
  \item \emph{Say it:} say \emph{lekhanam}
  \item \emph{Recall:} say the Sanskrit for a row, then the Sanskrit for a dictionary, then say \emph{lekhanam} again
\end{itemize}

\begin{tcolorbox}[breakable,colback=teal!4,colframe=teal!35!black,title={Before you move on}]
What does \sk{लेखनम्} mean? (\textquotedblleft{}Writing\textquotedblright{}.) Say it once more.
\end{tcolorbox}
